\section{Introduction}\label{sup:introduction}

A semistable limit cycle of a continuous piecewise smooth planar system may
have two spatially separated quadratic grazings, with both active sides
agreeing after transport to one transverse section. Each tangency adds a term
of order $q^{3/2}$ to the return map. When the leading coefficients of the two
terms cancel and the closing curvature is nonzero, the displacement has a
double zero at the reference parameter, but nearby return maps can have six
simple fixed points. The companion paper~\cite{main2026} determines the
cyclicity and the saturation domain at the principal scale. Here we study
smaller scales at a boundary of that domain.

The positive cluster at the principal scale is described by a quartic with
four nonnegative roots of sum one. When its quadratic coefficient tends to
zero, three roots approach zero and the fourth approaches one. Their next
coefficient region is cubic. A polynomial rescaling suggests this cubic, but
does not control the return map: the secondary scale can vanish at any rate
relative to the strength sum, and critical points can approach the moving
second threshold. The proof below uses divisions made before either scale
vanishes and slope estimates near that threshold.

We work with positive reference first passage strength $B$ and negative
reference closing curvature $-C$, where $B,C>0$. In terms of the offset
$a=R(0)$, the multiplier defect $b$, the strength sum $m$ and the threshold
separation $\delta$, the secondary weights are
\[
 m=\mu,\qquad b=-\mu^2\sigma,\qquad
 \delta=\mu^3\sigma^2d,\qquad a=\mu^4\sigma^3\alpha.
\]
Both $\mu$ and $\sigma$ are positive and independently small. The limiting
cubic is $z^3-z^2+(3Bd/2)z+\alpha$. Its coefficient region
$\mathcal C_B$ consists of cubics with three nonnegative roots of sum one.
Proposition~\ref{corner:complete-cluster-set} proves that this is exactly
the secondary cluster set and the Hausdorff limit of the six-cycle domains.
Necessity holds in every admissible family; realization uses the standard
unfolding of~\cite{main2026}.

Six-cycle configurations with positive offset occur on a finer offset scale.
With $h=-b$, Proposition~\ref{corner:secondary-bounds} proves the sharp
necessary bound $\limsup a/h^3\leq4/(27B^2)$.
Proposition~\ref{corner:positive-offset-realization} realizes every compact
interior set of the associated parameters by six simple cycles, one with
$q<0$ and five with $q>0$.

At the triple-root endpoint $(\alpha,d)=(-1/27,2/(9B))$,
Theorem~\ref{thm:actual-cusp} constructs an ordinary cusp for every
sufficiently small independent pair $(\mu,\sigma)$.
Lemma~\ref{corner:all-roots} locates every other root in the same interval
$I$. The three simple roots outside the cluster convert the local cubic counts
into the complete four/six-cycle diagram, including its folds and cusp.

The $q^{3/2}$ passage and the composition with regular transitions are
classical~\cite{dibbc2001,nordmark2002}; we use the root-profile form
established in~\cite[Section~2]{main2026}. The Malgrange
preparation theorem~\cite[Chapter~IV, Section~2, Theorem~2.1 and Note~(1), p.~95]{golubitsky1973} is applied after the
multiple root has been located away from both thresholds, where the
displacement is smooth in $q$.

Section~\ref{sup:setting} recalls the setting and the results
of~\cite{main2026} used here. Section~\ref{corner:section} passes from the
quartic corner to the cubic domain through divisions, necessary bounds and
root charts. Section~\ref{sup:positive} treats configurations with positive
offset. Section~\ref{sup:cusp} first locates all roots in $I$ and then unfolds
the secondary triple root.
